\section{Introduction}
\label{sec:intro}

Public football data matured over the last decade to the point where every outfield player
in a major league carried a rich statistical profile. Expected goals models
\citep{rathke2017examination, eggels2016explaining, vatvani2022upgrading} moved from
research novelty to broadcast furniture, and free sites published per-90 measures of
passing, carrying, pressing and chance creation that a decade earlier existed only inside
clubs. Most public analysis nonetheless stopped at leaderboards, ranking players on one
statistic at a time rather than asking what combinations of statistics describe a role.

Research on roles has generally required data that is not public. Bialkowski et al.
\citep{bialkowski2014large} used tracking data to recover formations and positional roles
from player coordinates. Decroos and Davis \citep{decroos2019player} learned player vectors
from event streams, Pappalardo et al. \citep{pappalardo2019playerank} built a performance
rating from event data across several leagues, and Aalbers and Van Haaren
\citep{aalbers2019distinguishing} distinguished playing styles from match events. All of
these depend on proprietary event or tracking feeds. The question of how much role
structure is recoverable from freely available season aggregates has received much less
attention, and the answer has just changed sharply.

On 20 January 2026 the operator of FBref \citep{fbref2026}, the most widely used free
source of advanced football statistics, removed every statistic derived from its
commercial provider after the licence was terminated and deletion was demanded. The
removal applied retroactively to completed seasons. Expected goals, expected assisted
goals, shot and goal creating actions, all touch, carry and take-on counts, all passing
volume and progression, and the entire advanced goalkeeping table ceased to be available.
What remains is close to the statistical vocabulary of two decades ago.

This paper asks what player role structure survives that withdrawal. We assemble the
statistics that remain, merge an independent expected goals source to restore the shot
quality block, and run a full role discovery pipeline: dimensionality reduction, clustering
with a selection rule fixed in advance, bootstrap stability, supervised validation,
goalkeeper and team level extensions, and a cross-season replication. Every figure and
every number in this paper is generated by code in an accompanying repository built on the
\texttt{soccerdata} library \citep{robberechts2026soccerdata}, and no value in the text was
transcribed by hand.

Five findings are worth stating in advance.

First, the withdrawal is a methodological trap and not merely a loss of columns. The
affected tables are still served with complete headers and one row per player, and only
the values are gone. A feature set assembled from the published schema resolves correctly
and produces an empty matrix, so feature selection must now be made against observed values
rather than documentation.

Second, correlations between statistics reverse sign across position groups. Of
\NPairsTotal{} feature pairs, \NPairsSignFlip{} change sign between defenders, midfielders
and forwards, and all ten of the most divergent pairs do. A pooled correlation matrix
therefore averages over opposing relationships and describes no actual population of
players.

Third, the surviving data supports a coarse binary division and nothing finer. The
selection rule returns two clusters in every scope, and the division it finds is attacking
against defensive involvement rather than a rediscovery of listed positions, with an
adjusted Rand index of only \ARIvsPosition{} against the position groups. A supervised
model reaches an accuracy of \LGBMAccuracy{} on the same features, so the weak agreement
reflects the clustering choosing a different axis rather than an uninformative feature set.

Fourth, a correction that public football analysis applies routinely is miscalibrated.
Defensive counts must be adjusted for how much time a team spends without the ball, and
the conventional adjustment assumes the two are directly proportional. Fitted against team
data that assumption fails: the elasticities are \ElasticityInterceptions{},
\ElasticityTackles{} and \ElasticityFouls{}, so the conventional correction overcorrects
by a factor of two to three and, applied here, replaces the confound it removes with a
larger one of the opposite sign.

Fifth, the global division is robust while the finer divisions inside position groups are
not. Bootstrap stability, cross-season replication and empirical Bayes shrinkage
together place the global partition at an adjusted Rand index of \ShrinkARIGlobal{} under
shrinkage toward the position mean, while the defender partition falls to \ShrinkARIDF{}
and fails all three tests. The archetypes that public data can still support are
therefore fewer, and coarser, than the ones it supported a year ago.
